\section{Experiments}
\label{sec:experiments}

\subsection{Experimental Setup}
\label{subsec:experimental_setup}

We use \texttt{personaplex-rl-seamless}~\citep{ohashi2026multifacetedinteractivityalignmentfullduplex}, a reinforcement-learning post-trained variant of PersonaPlex~\citep{11463413}, as the backbone for DiffuPlex.
We freeze the original backbone and train only the stream-specific MASK embeddings and future-gated LoRA parameters on Seamless Interaction~\citep{agrawal2025seamlessinteractiondyadicaudiovisual}, a large-scale dataset of face-to-face dyadic interactions.
Unless otherwise noted, we use a maximum lookahead of $W=16$, a confidence threshold of $\tau_{\mathrm{conf}}=0.80$, and a revision threshold of $\tau_{\mathrm{rev}}=-0.5$.
We evaluate two inference policies over the same masked-future predictor.
DiffuPlex-\textsc{Listen} consumes future positions only during sustained listening, with a minimum retained length of $m=6$, whereas DiffuPlex-\textsc{Speak} removes this state restriction and uses $m=2$.
The prediction at $h=1$ is always retained, ensuring that each backbone wake advances the interaction by at least one 80ms frame.
We use the original backbone sampling configuration and fix the generation seed to 0.
Further data processing, training, and inference details are provided in Appendices~\ref{asubsec:data_processing}--\ref{asubsec:inference}.

\subsection{Evaluation}
\label{subsec:evaluation}

We evaluate DiffuPlex along four axes: full-duplex interaction, general spoken-language capability, human evaluation, and inference efficiency.
For full-duplex interaction, we use Full-Duplex-Bench v1.0~\citep{11433838} to evaluate pause handling, backchanneling, turn-taking, and user interruption, and additionally evaluate on Full-Duplex-Bench v1.5~\citep{11463576} under overlapping and interfering speech.
We evaluate general spoken-language capability with VoiceBench~\citep{chen-etal-2026-voicebench} and OpenAudioBench~\citep{li2025baichuanaudiounifiedframeworkendtoend}, reporting both the assistant text stream ($s{\rightarrow}t$) and transcripts of generated speech ($s{\rightarrow}s$).
For human evaluation, we measure speech quality (NMOS) on the assistant response following an interruption and conversational quality (DMOS) over the full interaction, using user-interruption examples from Full-Duplex-Bench v1.0.
Finally, we measure inference efficiency on a fixed 60-clip subset of Full-Duplex-Bench v1.0, with 15 clips per task on a single NVIDIA RTX A6000 GPU.
We report wall-clock streaming, end-to-end GPU, and Core LM time, together with frames advanced per backbone wake and wake-frame latency.
Detailed evaluation, runtime, and human-evaluation protocols are provided in Appendices~\ref{asec:evaluation_details} and~\ref{asec:human_evaluation}.

\subsection{Baselines}
\label{subsec:baselines}

We compare against Moshi~\citep{defossez2024moshispeechtextfoundationmodel}, PersonaPlex~\citep{11463413}, and \texttt{personaplex-rl-seamless}, the autoregressive backbone used by DiffuPlex.
To test whether similar computational savings can be obtained simply by skipping some autoregressive backbone calls, we additionally construct three training-free deferred-execution baselines using the same frozen backbone: \textsc{PAD}, \textsc{Silero}, and \textsc{Smart Turn}.
On a deferred 80ms interaction step, these baselines do not run the backbone or predict the skipped frame; instead, they emit silence for the assistant and resume autoregressive generation at a later step.
\textsc{PAD} skips backbone calls while the assistant predicts silence, \textsc{Silero} while the user is speaking, and \textsc{Smart Turn} until a turn-ending detector predicts that the user has finished.
We cap the interval between autoregressive backbone calls at $R\in\{2,4,6,8\}$ interaction steps.
Further implementation details are in Appendix~\ref{asubsec:deferred_baselines}.
